
\فصل{ذینفعان}

\قسمت{تیم پروژه}

جدول~\ref{جدول:تیم} اعضای تیم و نقش هر یک در پروژه را نشان می‌دهد.

\begin{table}[ht]
\centering
\caption{اعضای تیم و مسئولیت‌ها}
\label{جدول:تیم}
\begin{tabular}{|>{\raggedleft\arraybackslash}p{3.5cm}|>{\raggedleft\arraybackslash}p{3cm}|>{\raggedleft\arraybackslash}p{6cm}|}
\hline
\rl{\textbf{نام}} & \rl{\textbf{نقش}} & \rl{\textbf{مسئولیت‌ها}} \\
\hline
\rl{ملیکا علیزاده} & \rl{Lead Team} & \rl{هماهنگی تیم، تسهیل جلسات، پیگیری پیشرفت، رفع موانع} \\
\hline
\rl{امیرحسین شهیدی} & \rl{Product Owner} & \rl{تعریف و اولویت‌بندی نیازمندی‌ها، نگهداری \lr{backlog}، تعامل با ذینفعان} \\
\hline
\rl{معین آعلی} & \rl{Architecture Owner} & \rl{طراحی معماری سامانه، تصمیم‌های فنی کلان، نظارت بر کیفیت فنی کد} \\
\hline
\rl{محمدمهدی فراهانی} & \rl{Developer} & \rl{پیاده‌سازی کدبیس، نوشتن تست، ارائه‌ی راهکارهای فنی} \\
\hline
\end{tabular}
\end{table}

\قسمت{ذینفعان مستقیم}

ذینفعان مستقیم کسانی هستند که مستقیماً از ساکنا استفاده می‌کنند:

\زیرقسمت{ساکنین}
ساکنین واحدهای مسکونی اصلی‌ترین کاربران سامانه هستند.
آن‌ها نیاز دارند بدون مراجعه‌ی حضوری شارژ پرداخت کنند، درخواست تعمیرات ثبت نمایند،
امکانات مشترک رزرو کنند و اطلاعیه‌های ساختمان را دنبال کنند.
\textbf{انتظار:} رابط کاربری ساده، دسترسی ۲۴/۷ و پاسخ‌گویی سریع سامانه.

\زیرقسمت{مدیر ساختمان}
مدیر ساختمان مسئول اداره‌ی امور مجتمع است.
او نیاز دارد اطلاعات واحدها و ساکنین را مدیریت کند، درخواست‌های خدماتی را ارجاع دهد،
شارژها را تعریف کند و گزارش‌های مالی و عملیاتی دریافت نماید.
\textbf{انتظار:} داشبورد جامع، گزارش‌های دقیق و ابزار مدیریت ساده.

\زیرقسمت{کارکنان خدماتی}
کارکنان خدماتی (تکنسین‌ها و سرویس‌کارها) مسئول انجام تعمیرات و خدمات هستند.
آن‌ها نیاز دارند وظایف ارجاع‌شده را مشاهده کنند، وضعیت انجام کار را ثبت نمایند
و تاریخچه‌ی فعالیت‌های خود را ببینند.
\textbf{انتظار:} رابط ساده، لیست واضح وظایف و امکان گزارش‌دهی آسان.

\قسمت{ذینفعان غیرمستقیم}

\زیرقسمت{مدیر سامانه}
مدیر سامانه (\lr{System Admin}) مسئول نگهداری و مدیریت کلی پلتفرم است.
او می‌تواند فهرست کاربران را مشاهده کند و در صورت لزوم حساب‌های مشکل‌دار را غیرفعال نماید.

\زیرقسمت{استاد درس}
دکتر رامان رامسین به‌عنوان استاد درس «ایجاد چابک نرم‌افزار» ذینفع غیرمستقیم پروژه است.
موفقیت پروژه از دیدگاه آموزشی شامل یادگیری متدولوژی \lr{DAD}، اجرای صحیح \lr{Lean DAD}
و انتقال به \lr{Continuous Delivery Lean} در فاز دوم است.

\زیرقسمت{دستیاران آموزشی (\lr{TA})}
دستیاران آموزشی درس «ایجاد چابک نرم‌افزار» نیز ذینفعان غیرمستقیم پروژه هستند.
آن‌ها مسئول ارزیابی مستندات، بررسی تحویل‌دادنی‌ها در پایان هر ایتریشن
و ارائه‌ی بازخورد فنی و فرآیندی به تیم هستند.
\textbf{انتظار:} مستندات کامل و به‌روز، رعایت دقیق متدولوژی \lr{DAD}، شفافیت در \lr{backlog} و گزارش پیشرفت.


\زیرقسمت{تیم ایجاد ساکنا}
تیم ایجاد ساکنا مسئول تحلیل، طراحی، توسعه و تحویل سامانه است.
اعضای تیم با بهره‌گیری از متدولوژی \lr{DAD} و رویکردهای چابک، نیازمندی‌های ذینفعان را
به قابلیت‌های قابل استفاده در سامانه تبدیل می‌کنند. همچنین تیم وظیفه دارد کیفیت فنی،
پایداری سامانه، مستندسازی مناسب و تحویل به‌موقع خروجی‌ها را تضمین کند.
اعضای تیم به‌صورت مستمر با مدیر ساختمان، ساکنین و سایر ذینفعان در ارتباط هستند تا
بازخوردهای دریافتی را در فرآیند توسعه لحاظ کنند و اطمینان حاصل شود که محصول نهایی
نیازهای واقعی کاربران را برآورده می‌کند.

